\section{Challenges in Evaluating \gai}
\label{section:13-graph-based-ai-challenges}


\begin{figure}[h!]
    \begin{center}
        \includesvg[width=\columnwidth]{SVGs/gai_challenges}
        \caption{\gai challenges are a hierarchy of dependencies foundationally limited by graph processing.}
        \label{figure:gai-challenges}
    \end{center}
\end{figure}



A hierarchy of challenges affects \gai (Figure~\ref{figure:gai-challenges}). 
This section briefly explains the core problems at each level of the hierarchy and shows how this thesis's research aim is grounded in addressing these challenges. 

\textbf{Graph Challenges.} All problems in the hierarchy build on the foundational challenge of data dependence in graph processing
Graphs of different shapes and sizes vary the algorithmic complexity and exacerbate the irregular memory accesses and computation patterns that undermine spatial and temporal locality. 
Graphs have no standard accepted representation method, abstractions, or primitive operations.
The resulting landscape consists of competing approaches and no unity of effort in accelerating graph processing generally. 

\textbf{\gai Challenges.} Graph processing is the bottleneck in \gai, meaning that improving \gai performance is directly dependent on elevating the graph processing bottlenecks. 
It is difficult to achieve a universal improvement in \gai because each problem in the taxonomy has different computational characteristics and behaves differently on the same data. 


\textbf{Real-World Challenges.} Graph and \gai processing difficulties are exacerbated by the constraints inherent to deploying them in the real world. 
In particular, \gai solutions must be delivered within the efficiency, effectiveness, and cost constraints. 
For example, where a real-world problem requires an hourly solution but the \gai task that generates it requires ninety minutes to run, compromises must be made to the effectiveness or cost to meet the efficiency constraint. 
Addressing this problem is not as simple as just `accelerating' the \gai task.
For example, hardware accelerators for graph processing are specialized \ac{ASIC} optimized for the specific algorithm, representation, abstraction, and programming model.
Because one system may need to support many \gai tasks, an accelerator will need to improve the performance of all those tasks to be a wise investment. 

\textbf{User Challenges.} The user then must decide whether to trade effectiveness for efficiency or pay a higher cost to achieve the desired results.
Gathering the data to make that decision has historically included choosing between a simple competitive benchmark of systems or an invasive workload analysis to determine which combination of hardware, representation, abstraction model, and algorithm will best meet requirements. 

\textbf{The Evaluation Challenge.}The challenges culminate in the research contribution of this thesis, \textit{determining \textit{if}, \textit{how} and \textit{when} to speed up problem-solving on linked data given a set of operating constraints.}